\section*{الفصل \ref{ch:g4:magnets-and-compass} --- المغانط والبوصلة}

\begin{solution}{exo:g4:magnets-and-compass:1}
القطبان هما طرفا \omterm{def:g1:magnets:magnet}{المغناطيس} القويان: \omterm{def:g4:magnets-and-compass:poles}{القطب الشمالي} (N، وكثيرًا ما
يكون أحمر) \omterm{def:g4:magnets-and-compass:poles}{والقطب الجنوبي} (S). وتتجمّع مشابك الورق عند القطبين
--- أي الطرفين --- لا عند الوسط.
\end{solution}

\begin{solution}{exo:g4:magnets-and-compass:2}
القطبان المختلفان يتجاذبان؛ والمتماثلان يتنافران. والالتصاق
يعني أن قطبين مختلفين كانا متقابلين: قطب N في مواجهة قطب S.
\end{solution}

\begin{solution}{exo:g4:magnets-and-compass:3}
N نحو N: يتدافعان (فالمتماثلان يتنافران). وN نحو S: يتجاذبان
ويلتصقان.
\end{solution}

\begin{solution}{exo:g4:magnets-and-compass:4}
لأن مغناطيسًا عملاقًا يسحب قطبيه إلى الاصطفاف: فالأرض نفسها
\omterm{def:g1:magnets:magnet}{مغناطيس}، \omterm{def:g4:magnets-and-compass:poles}{وقطباه المغناطيسيان} قرب أعلى الكرة الأرضية وأسفلها.
\end{solution}

\begin{solution}{exo:g4:magnets-and-compass:5}
الشمال، فالشرق، فالجنوب، فالغرب. وأنت مواجه الشمال يكون الشرق
عن يمينك.
\end{solution}

\begin{solution}{exo:g4:magnets-and-compass:6}
الإبرة \omterm{def:g1:magnets:magnet}{مغناطيس} دقيق يطيع الجذب \emph{الأقرب} الأقوى. \omterm{def:g1:magnets:magnet}{ومغناطيس}
أو \omterm{def:g3:measuring-mass:mass}{كتلة} حديد كبيرة قريبة تغلب \omterm{def:g1:magnets:magnet}{مغناطيس} الأرض البعيد، فتكذب
الإبرة في شمالها.
\end{solution}

\begin{solution}{exo:g4:magnets-and-compass:7}
دلّك الإبرة بأحد قطبَي \omterm{def:g1:magnets:magnet}{مغناطيس} عشرين مرة في الجهة نفسها، مبعدًا
\omterm{def:g1:magnets:magnet}{المغناطيس} بين الدلكة والأخرى --- فتصير مغناطيسًا صغيرًا. ثم
ضعها على شريحة فلّين تطفو في وعاء ماء: فالفلّينة الطافية \omterm{def:g3:levers-and-scales:lever}{محور
ارتكاز} من ماء، وتدور الإبرة حتى تشير شمالًا وجنوبًا.
\end{solution}

\begin{solution}{exo:g4:magnets-and-compass:8}
يشير إلى الشمال. \omterm{def:g4:magnets-and-compass:compass}{والبوصلة} التي توافق إبرتُها ظلَّ الظهر تقول
الصدق --- وفي يوم مشمس يصلح \omterm{def:g1:light-and-shadows:shadow}{الظلّ} وحده \omterm{def:g4:magnets-and-compass:compass}{بوصلة} تقريبية.
\end{solution}

\begin{solution}{exo:g4:magnets-and-compass:9}
لأنها تعمل في الضباب، وتحت الغيوم، وفي الليالي بلا \omterm{def:g2:moon-in-the-sky:moon}{قمر}، وبعيدًا
عن أي ساحل --- فالإبرة لا تطلب من السماء شيئًا. فاستطاع
البحّارة أخيرًا أن يحفظوا مسارًا عبر المحيطات المفتوحة في كل
طقس.
\end{solution}

\begin{solution}{exo:g4:magnets-and-compass:10}
إبزيم صديقها الحديدي يغلب \omterm{def:g1:magnets:magnet}{مغناطيس} الأرض عن قرب: فتشير الإبرة
إلى الإبزيم لا إلى الشمال. ولو وثقت بها لانحرفت --- فالشمال
الكاذب يقع نحو يسارها، فتنثني بمسارها يسارًا، أي غربًا. والحلّ:
أن تقرأ \omterm{def:g4:magnets-and-compass:compass}{البوصلة} بعيدًا عن الإبزيم (ببضع خطوات فاصلة)، ثم
تمشي.
\end{solution}

\begin{solution}{exo:g4:magnets-and-compass:11}
الإبرة القصيرة المدلوكة لها قطبان \emph{اثنان} --- فطرفاها
يتصرّفان تصرّف N وS (أطفِ إبرتين كهاتين فتلتصقان طرفًا إلى
طرف، ويتجاذب الطرفان المختلفان). إذن تقصير \omterm{def:g1:magnets:magnet}{المغناطيس} لا ينزع
منه قطبًا، وقطعه نصفين لا يصنع إلا مغناطيسين قصيرين كاملين:
فاصطياد قطب وحده بالقطع لا ينتهي أبدًا --- إذ يسلّمك كل قطع
القطبين من جديد.
\end{solution}
